\section{Introduction}

Visual generation papers compare models through finite-sample metric estimates:
a single FID, KID, or CMMD number at some chosen sample size declares a winner.
A point estimate alone does not say whether one model is \emph{certifiably}
better under a declared metric and budget, and the way these numbers are
produced invites trouble. Evaluators choose the sample size, the seeds, the
checkpoint, and which metric to report, and a comparison is often read as decided
the moment its gap looks large. This is a garden of forking paths: fixed-sample
confidence intervals are not valid under data-dependent stopping and selection,
and a small reported gap may be statistical noise rather than a real difference.

CertGen frames model comparison as a decision problem. Given released samples or
reproducible feature caches, a reference set, a preprocessing lock, and a metric
stream, it asks whether the optional-stopping-safe uncertainty interval for the
paired discrepancy $\Delta_{AB}=d(A,R)-d(B,R)$ excludes zero before the sample
budget is exhausted. Validity is provided by a betting confidence sequence on the
bounded per-block stream, so the decision survives continuous monitoring; the
same e-process yields an e-value per comparison, and e-BH controls the false
discovery rate across an entire leaderboard under the arbitrary dependence
induced by a shared reference set. Two outputs follow naturally: a
\emph{samples-to-decision} number (how cheaply a real win can be certified, i.e.,
a compute-savings story for evaluation) and a leaderboard \emph{audit} (what
fraction of reported wins are statistically decided at the budgets actually
used).

CertGen is a decision layer, not a new metric, and the statistics are borrowed
rather than invented. It separates clean-core metrics with bounded streamable
estimators (MMD, CMMD) from descriptive metrics such as FID and FD-DINOv2, and it
gates every claim behind a provenance ledger, a preprocessing lock, and a
reproduction check so that an audit on real benchmarks is defensible. The
empirical audit sentence is reserved for claim-eligible real runs:
TBD\_REAL\_RUN\_REQUIRED.
